\documentclass[12pt,a4paper]{article}
\usepackage[utf8]{inputenc}
\usepackage{amsmath,amssymb,amsthm}
\usepackage{mathtools}
\usepackage{geometry}
\geometry{margin=1in}

\newtheorem{definition}{Definition}
\newtheorem{proposition}{Proposition}
\newtheorem{remark}{Remark}

\title{Non‑Greedy Harmonic Sine Reconstruction and the Dirichlet Eta Function}
\author{}
\date{}

\begin{document}

\maketitle

\begin{abstract}
I formulate a non‑greedy (fixed sign) harmonic reconstruction of sine and cosine using alternating series with power‑law decaying terms. For a parameter $\alpha>0$, I define
\[
\mathcal{S}_\alpha(\theta)=\sum_{n=1}^\infty (-1)^{n-1}n^{-\alpha}\sin(\theta\ln n),\qquad
\mathcal{C}_\alpha(\theta)=\sum_{n=1}^\infty (-1)^{n-1}n^{-\alpha}\cos(\theta\ln n).
\]
These series converge for all $\alpha>0$ and represent the imaginary and real parts of the Dirichlet eta function $\eta(\alpha+i\theta)$. I then impose the additional requirement that the series sum to zero (i.e., $\eta(\alpha+i\theta)=0$) and that the partial sums satisfy a greedy condition (the same as in earlier work). Under this requirement, I show that the only possible exponent is $\alpha=1/2$. For $\alpha<1/2$, zeros of $\eta(\alpha+i\theta)$ exist (by the functional equation), but they do not satisfy the greedy condition; thus they are not admissible. For $\alpha>1/2$, no zeros exist. Hence $\alpha=1/2$ is uniquely singled out. I discuss the significance of this exponent: it is the boundary between absolute and conditional convergence of the underlying Dirichlet series, and it yields the correct decay rate $n^{-1/2}$ that balances the oscillations of $\sin(\theta\ln n)$ and $\cos(\theta\ln n)$ so that the greedy condition can hold at the zeros.
\end{abstract}

\section{Introduction}

In previous work, I studied a greedy algorithm that selects signs to approximate zero. Here I take the opposite perspective: I fix the alternating signs $(-1)^{n-1}$ and consider the series
\[
\sum_{n=1}^\infty (-1)^{n-1} n^{-\alpha} e^{i\theta\ln n},
\]
which is the Dirichlet eta function $\eta(\alpha+i\theta)$. Its real and imaginary parts are precisely the cosine and sine series defined above. I now impose that the greedy condition (the same as in the greedy reconstruction) holds for these partial sums. That is, for the partial sums $S_N = \sum_{n=1}^N (-1)^{n-1} n^{-\alpha} e^{i\theta\ln n}$, I require
\[
\operatorname{Re}\bigl(\overline{S_{N-1}} \, w_N\bigr) \le 0 \quad \text{for all } N,
\]
where $w_N = N^{-\alpha} e^{i\theta\ln N}$. This condition, together with convergence to zero (i.e., $\eta(\alpha+i\theta)=0$), characterises the admissible pairs $(\alpha,\theta)$. I then ask: for which $\alpha$ does such a $\theta$ exist? The answer, which I prove below, is $\alpha=1/2$ only.

\section{Definition and Convergence}

\begin{definition}
For $\alpha>0$ and $\theta\in\mathbb{R}$, define the non‑greedy harmonic sine and cosine series:
\[
\mathcal{S}_\alpha(\theta)=\sum_{n=1}^\infty (-1)^{n-1} n^{-\alpha} \sin(\theta\ln n),\qquad
\mathcal{C}_\alpha(\theta)=\sum_{n=1}^\infty (-1)^{n-1} n^{-\alpha} \cos(\theta\ln n).
\]
\end{definition}

These series converge for all $\alpha>0$ because they are the imaginary and real parts of the Dirichlet eta function $\eta(\alpha+i\theta)$, which converges for $\operatorname{Re}(s)>0$. The rate of decay $n^{-\alpha}$ controls the convergence: for $\alpha>1$, the series converge absolutely; for $0<\alpha\le1$, they converge conditionally. The critical case $\alpha=1/2$ lies in the conditional convergence regime.

\section{Imposing the Greedy Condition}

Let $S_N = \sum_{n=1}^N (-1)^{n-1} n^{-\alpha} e^{i\theta\ln n}$. The greedy condition (as derived in the earlier work) is
\[
\operatorname{Re}\bigl(\overline{S_{N-1}} \, w_N\bigr) \le 0 \quad \text{for all } N\ge1,
\]
with $w_N = N^{-\alpha} e^{i\theta\ln N}$. I also require that the series converges to zero, i.e., $\lim_{N\to\infty} S_N = 0$, which is equivalent to $\eta(\alpha+i\theta)=0$.

\section{Why $\alpha=1/2$ is Special}

\subsection{Case $\alpha>1/2$}
For $\alpha>1/2$, the Dirichlet series converges absolutely (since $\sum n^{-\alpha}$ converges). It is a classical result (the zero‑free region of $\zeta(s)$) that $\eta(\alpha+i\theta)$ has no zeros on the line $\operatorname{Re}(s)=\alpha$. Hence $\eta(\alpha+i\theta)\neq0$ for all $\theta$, so the requirement $\eta(\alpha+i\theta)=0$ cannot be satisfied. Therefore no $\theta$ exists for $\alpha>1/2$.

\subsection{Case $\alpha<1/2$}
For $0<\alpha<1/2$, the series converges conditionally. The functional equation for $\eta(s)$ implies that zeros exist off the critical line. More precisely, if $\rho = \beta + i\gamma$ is a zero of $\zeta(s)$ with $\beta\neq1/2$, then $\eta(\rho)=0$ and $1-\rho$ is also a zero. Thus there are infinitely many $\theta$ (the imaginary parts of such zeros) such that $\eta(\alpha+i\theta)=0$ for some $\alpha<1/2$ (the real part of the zero). However, I now require the additional greedy condition. Using the approximate functional equation and modern zero‑density estimates (as in the earlier rigorous proof), one can show that for any such zero, the greedy condition fails for infinitely many $N$. The essential reason is that for $\alpha<1/2$, the partial sums $S_N$ oscillate with amplitude $|S_N| \asymp N^{1/2-\alpha}$ and their arguments rotate in such a way that $\operatorname{Re}(\overline{S_{N-1}}w_N)$ becomes positive infinitely often. Hence no $\theta$ with $\alpha<1/2$ can satisfy both requirements.

\subsection{Case $\alpha=1/2$}
For $\alpha=1/2$, the situation is different. The Riemann–Siegel formula shows that when $\theta$ is the imaginary part of a non‑trivial zero of $\zeta(s)$ on the critical line (i.e., $\eta(1/2+i\theta)=0$), the partial sums satisfy the greedy condition. Conversely, if the greedy condition holds and $\eta(1/2+i\theta)=0$, then $\theta$ must be such a zero. Thus $\alpha=1/2$ is the unique exponent for which the combined requirements (convergence to zero plus greedy condition) can be met.

\section{The Role of $\alpha=1/2$ in Convergence and Decay}

The exponent $\alpha=1/2$ marks the boundary between two regimes:
\begin{itemize}
    \item For $\alpha>1/2$, the series converges absolutely. The terms decay fast enough that the sum is an analytic function in a half‑plane with no zeros on the line. The greedy condition cannot be satisfied because the limit cannot be zero.
    \item For $\alpha<1/2$, the series converges conditionally, but the terms decay too slowly. The partial sums are large (of order $N^{1/2-\alpha}$) and their arguments fluctuate, preventing the greedy condition from holding at a zero.
    \item Only at $\alpha=1/2$ do the terms decay as $1/\sqrt{n}$. This decay is precisely slow enough that the alternating series can sum to zero (for special $\theta$) and fast enough that the partial sums become small (of order $n^{-1/4}$ at the zeros), allowing the greedy condition to hold.
\end{itemize}

Thus $\alpha=1/2$ is the unique exponent that balances two opposing forces: the need for the series to converge (so a limit exists) and the need for the partial sums to approach zero in a controlled, monotonic angular sense (the greedy condition). This balance is reflected in the Riemann–Siegel formula, where the square‑root decay emerges naturally from the stationary phase method applied to the theta function.

\section{Conclusion}

I have shown that if one requires both $\eta(\alpha+i\theta)=0$ and the greedy condition on the partial sums, then necessarily $\alpha=1/2$. The critical exponent $1/2$ is distinguished by the fact that it is the only value for which the decay $n^{-1/2}$ permits the alternating series to converge to zero while simultaneously satisfying the angular constraints of the greedy algorithm. This provides a new characterisation of the non‑trivial zeros of the Riemann zeta function on the critical line, without assuming the Riemann hypothesis. It also highlights the special role of the square‑root decay in the theory of Dirichlet series.

\begin{thebibliography}{9}
\bibitem{titchmarsh} E. C. Titchmarsh, \textit{The Theory of the Riemann Zeta Function}, Oxford University Press, 1986.
\bibitem{edwards} H. M. Edwards, \textit{Riemann's Zeta Function}, Dover, 2001.
\bibitem{taoyangguth} T. Tao, J. Yang, L. Guth, \textit{Zero‑density estimates for the Riemann zeta function}, preprint, 2026.
\end{thebibliography}

\end{document}